\clearpage
\VerdierPhysicalPageAnchor{0143}%
\begingroup
\small
\setlength{\parskip}{0pt}%
\setlength{\emergencystretch}{1.5em}%

\begingroup\itshape
\VerdierSemanticAnchor{ch2-2-3-1-continuation}%
For every full triangulated subcategory $\mathcal A$ such that
$\mathcal B\subset\mathcal A\subset\mathcal D$, denote by
$S_{\mathcal D}(\mathcal B)$ (respectively $S_{\mathcal A}(\mathcal B)$)
the multiplicative system of $\mathcal D$ (respectively $\mathcal A$)
associated with $\mathcal B$ (2.1.7).

\VerdierSemanticAnchor{ch2-2-3-1-a-intersection}
\par\noindent a) One has
$S_{\mathcal D}(\mathcal B)\cap\operatorname{Fl}(\mathcal A)
=S_{\mathcal A}(\mathcal B)$.

\VerdierSemanticAnchor{ch2-2-3-1-b-fully-faithful}
\par\noindent b) The canonical functor
$\mathcal A/\mathcal B\longrightarrow\mathcal D/\mathcal B$ is fully
faithful and injective on objects. The image of this functor is
$Q(\mathcal A)$, where
$Q\colon\mathcal D\longrightarrow\mathcal D/\mathcal B$ is the canonical
functor.

\VerdierSemanticAnchor{ch2-2-3-1-c-quotient-transitivity}
\par\noindent c) The canonical functor
$\mathcal D/\mathcal A\longrightarrow
(\mathcal D/\mathcal B)/(\mathcal A/\mathcal B)$ is an isomorphism of
categories.

\VerdierSemanticAnchor{ch2-2-3-1-d-saturation}
\par\noindent d) If the category $\mathcal A$ is saturated (2.1.6), the
subcategory $Q(\mathcal A)$ is saturated.

\VerdierSemanticAnchor{ch2-2-3-1-inverse-image-setup}
For every full triangulated subcategory $\mathcal C$ of
$\mathcal D/\mathcal B$, denote by $Q^{-1}(\mathcal C)$ the full
subcategory of $\mathcal D$ defined by the objects $X$ of $\mathcal D$
such that $Q(X)$ is isomorphic to an object of the category $\mathcal C$.

\VerdierSemanticAnchor{ch2-2-3-1-a-bis}
\par\noindent a)$^{\mathrm{bis}}$ The category $Q^{-1}(\mathcal C)$ is a
strictly full triangulated subcategory of $\mathcal D$.

\VerdierSemanticAnchor{ch2-2-3-1-b-bis}
\par\noindent b)$^{\mathrm{bis}}$ The multiplicative system
$S_{\mathcal D}(Q^{-1}(\mathcal C))$ consists of the morphisms of
$\mathcal D$ whose image under $Q$ is a morphism in
$S_{\mathcal D/\mathcal B}(\mathcal C)$.

\VerdierSemanticAnchor{ch2-2-3-1-c-bis}
\par\noindent c)$^{\mathrm{bis}}$ The subcategory $Q^{-1}(\mathcal C)$ is
saturated if and only if the subcategory $\mathcal C$ is saturated.

\VerdierSemanticAnchor{ch2-2-3-1-d-bis}
\par\noindent d)$^{\mathrm{bis}}$ The map
$\mathcal A\longmapsto Q(\mathcal A)$ is an isomorphism of ordered sets
from the set of strictly full triangulated subcategories of $\mathcal D$
containing $\widehat{\mathcal B}$ (the kernel of the functor $Q$), ordered
by inclusion, to the set of strictly full triangulated subcategories of
$\mathcal D/\mathcal B$, ordered by inclusion. The inverse map is
$\mathcal C\longmapsto Q^{-1}(\mathcal C)$.
\endgroup

\VerdierSemanticAnchor{ch2-2-3-1-proof-ab}
The first assertion is evident. The second follows immediately from the
description of morphisms in the categories $\mathcal A/\mathcal B$ and
$\mathcal D/\mathcal B$ (2.2.1).

\VerdierSemanticAnchor{ch2-2-3-1-proof-c}
To prove c), it suffices to refer to the universal properties of quotient
categories (2.2.11).

\VerdierSemanticAnchor{ch2-2-3-1-proof-d}
Let us prove d). After replacing $\mathcal A$ by its strictly full closure,
which changes $Q(\mathcal A)$ only up to equivalence, we may assume that
$\mathcal A$ is strictly full. The kernel of the functor
$\mathcal D\longrightarrow\mathcal D/\mathcal A$ is then the category
$\mathcal A$ (2.2.11). This functor factors as
$\mathcal D\longrightarrow\mathcal D/\mathcal B
\longrightarrow\mathcal D/\mathcal A$. It follows that the kernel of the
canonical functor
$\mathcal D/\mathcal B\longrightarrow\mathcal D/\mathcal A$ is
$Q(\mathcal A)$. Consequently, $Q(\mathcal A)$ is saturated (2.1.16).

\VerdierSemanticAnchor{ch2-2-3-1-proof-a-bis-b-bis}
Assertion a)$^{\mathrm{bis}}$ is evident, as is assertion
b)$^{\mathrm{bis}}$.

\VerdierSemanticAnchor{ch2-2-3-1-proof-c-bis-forward}
Suppose that $\mathcal C$ is saturated. We may further assume that
$\mathcal C$ is strictly full, for if $\widehat{\mathcal C}$ denotes the
strictly full closure of $\mathcal C$, then the category
$\widehat{\mathcal C}$ is strictly full and saturated and
$Q^{-1}(\widehat{\mathcal C})=Q^{-1}(\mathcal C)$. The category
$\mathcal C$ is then the kernel of an exact functor (2.2.11).
Consequently, $Q^{-1}(\mathcal C)$ is the kernel of an exact functor, hence
is saturated (2.1.16).

\VerdierSemanticAnchor{ch2-2-3-1-p144-boundary}
Suppose that $Q^{-1}(\mathcal C)$ is
\endgroup
